% ============================================================================
\section{Discussion and Conclusion}
\label{sec:conclusion}
% ============================================================================

GPAS treats task importance and gradient-estimation effort as separate
decisions. Existing dense distillation provides the learning signal;
the scheduler uses its observed variability to decide where more
micro-batches are useful. The same principle gives a per-step rule and a
cost-aware rule, both using statistics collected during ordinary training.
\missing{Conclude with the measured efficiency difference and the task-level
learning pattern that explains it.}

The present study covers a Qwen3-1.7B student, four domains, one training
seed per method, and one hardware layout. The method is most directly
applicable when an optimizer step contains several micro-batches from each
task; larger task collections call for less frequent task measurements.
The pre-step optimizer scaling and measured time model provide local
allocation signals, and their practical effect is evaluated with the
complete AdamW training loop. Task compatibility and student capacity
remain important for capability integration. GPAS offers a way to improve
how a chosen collection of task objectives uses computation.
